% Copyright (c) 2026 Eugene Joseph M. Ragasa.
% SPDX-License-Identifier: Apache-2.0
% Author: Eugene Joseph M. Ragasa
% Formal proof completed: 2026-10-27
\documentclass[11pt]{article}
\usepackage[T1]{fontenc}
\usepackage{amsmath,amssymb}
\usepackage[margin=1in]{geometry}
\usepackage[hidelinks]{hyperref}
\title{Proof Dossier: Global-Unitary Feasibility}
\author{Eugene Joseph M. Ragasa}
\date{Formal proof completed: 2026-10-27}
\begin{document}
\raggedright
\maketitle

\section{Identity and status}
Statement ID \texttt{BG-ALIGNMENT-THM-018}; status \texttt{proved; no admissions}.  Lean declaration
\path{ProjectKoios.BandGauge.Alignment.alignmentFeasible_global_iff} in
\path{ProjectKoios/BandGauge/Alignment/Family.lean}.

\section{Exact mathematical statement}
Feasibility in the global family is equivalent to existence of one unitary $U$ whose constant alignment field is accepted.

\section{Assumptions and conventions}
One common unitary must act at every point.  The theorem allows different reference and candidate fiber types when a linear isometric equivalence exists.

\section{Mathematical argument}
A global-family membership proof already contains a unitary $U$ and equality with its constant field.  Rewriting the feasible witness by that equality gives the forward existential.  Conversely, an accepted constant field is paired with $U$ as its global-family membership witness.

\section{Lean realization}
The Lean proof follows the mathematical argument directly using the definitions and the
previously proved helper declarations in the same module family.  The declaration was checked
with Lean \texttt{4.35.0} and Mathlib commit \texttt{56a005855ea17e0b0958f73ec801ae140b0a1679}.

\section{Provenance and contribution boundary}
This is source-neutral finite mathematics at Formal merge
\texttt{d9fe61eba430fd2c101dc53d85a2e0978254a718}.  No external manuscript controls the Lean statement, and no priority or
novelty claim is made for the underlying algebra or set theory.

\section{Audit result}
The named axiom report lists \texttt{propext}, \texttt{Classical.choice}, and
\texttt{Quot.sound} and no repository-owned axiom.  The repository admission audit found no
\texttt{sorry} or \texttt{admit} in the checked source.

\section{Boundary and nonclaims}
The unitary may be existentially nonunique, but it is common to every point.  This theorem does not permit momentum-dependent alignment.

\section{Review status}
AI-assisted exposition, checked against the Lean source and theorem manifest.  Kernel checking
establishes the exact proposition only.  No independent expert mathematical, scientific,
novelty, or publication review is recorded by this dossier.

\end{document}
